% Libro 3 · Mundo 2 · Tipos de texto y géneros (app: pu-tipos)
\mundo{Tipos de texto y géneros}{pu-tipos}{No se lee igual un cuento que un editorial. Reconocer qué clase de texto se
tiene enfrente dice qué buscar: en un relato, quién hace qué y qué cambia; en un ensayo, la tesis y sus razones; en un
texto expositivo, la información y cómo está organizada. El examen pregunta directamente por el tipo de texto y, sobre
todo, lo usa para todo lo demás.}{%
\item \textbf{2.1} Tipologías: narrar, describir, explicar, argumentar
\item \textbf{2.2} Los géneros que aparecen en el examen
\item \textbf{2.3} Estructura del texto y testimonio}

% ------------------------------------------------------------------
\leccion{pu-tt-1}{Tipologías}{Juego de clasificar fragmentos por su intención.}

\begin{reto}
«Para preparar el suero casero, disuelva en un litro de agua hervida y fría seis cucharaditas rasas de azúcar y media
cucharadita de sal. Revuelva bien y consérvelo tapado.» Este fragmento es principalmente
\opcionesd{narrativo}{argumentativo}{instructivo}{descriptivo}
\end{reto}

\tecnica{La técnica: preguntarle al texto qué quiere hacer}
Cada \clave{tipología} responde a una intención distinta. La forma más rápida de reconocerla es preguntar: ¿qué quiere
lograr el autor con esto?

\begin{center}\small
\begin{tabularx}{\linewidth}{@{}l>{\raggedright\arraybackslash}X>{\raggedright\arraybackslash}X@{}}
\toprule
\textbf{Tipo} & \textbf{Intención} & \textbf{Marcas}\\\midrule
Narrativo & contar hechos en el tiempo & personajes, verbos en pasado, «después», «de pronto»\\
Descriptivo & mostrar cómo es algo & adjetivos, comparaciones, verbos de estado\\
Expositivo & explicar, informar & definiciones, datos, ejemplos, tono neutro\\
Argumentativo & convencer de una postura & tesis, razones, «por lo tanto», «sin embargo»\\
Instructivo & decir cómo hacer algo & imperativos, pasos numerados\\
Dialogal & intercambio entre voces & rayas, preguntas y respuestas\\\bottomrule
\end{tabularx}
\end{center}

\subsubsection{Textos mezclados}
Casi ningún texto real es puro: una crónica narra, describe y a veces opina. La pregunta del examen se refiere a la
tipología \emph{predominante}, la que organiza el texto.

\figura{tipologias}{Cinco intenciones principales y sus tipos de texto.}

\begin{regla}
\begin{itemize}
  \item Pregunta «¿para qué?» y no «¿de qué?»: un texto sobre el agua puede contar, explicar o convencer.
  \item Busca las marcas: imperativos (instructivo), tesis y conectores de razonamiento (argumentativo), datos y
  definiciones (expositivo), secuencia temporal (narrativo).
  \item Elige la tipología que domina, no la que aparece en una frase.
\end{itemize}
\end{regla}

\begin{texto}{Cuatro fragmentos sobre la lluvia}{Fragmentos originales escritos para este libro}
\pa La lluvia comenzó a las tres. Rosa recogió la ropa del patio, cerró las ventanas y se sentó a esperar a su hijo, que
había salido en bicicleta. A las cinco, empapado y feliz, Tomás tocó la puerta.

\pa La llovizna de los páramos es fina como polvo y no cae: flota. Moja sin que uno lo note, deja en la lana de las ruanas
un brillo de vidrio molido y vuelve gris todo lo que toca.

\pa La lluvia se forma cuando el vapor de agua de la atmósfera se enfría y se condensa en gotitas alrededor de partículas
de polvo. Cuando las gotitas se unen y alcanzan cierto tamaño, el aire ya no puede sostenerlas y caen.

\pa Las ciudades deberían recoger el agua de lluvia de sus techos. Es gratuita, reduce la presión sobre los acueductos y
disminuye las inundaciones. Que se siga botando por las canaletas es un desperdicio que ya no podemos permitirnos.
\end{texto}

\subsection{Preguntas comentadas}
\begin{comentada}{El fragmento del \S\,4 es predominantemente argumentativo porque}
\mal{A}{habla de la lluvia.}{El tema no define la tipología: los cuatro fragmentos hablan de la lluvia.}
\mal{B}{da datos científicos.}{\emph{Dice otra cosa}: eso caracteriza al \S\,3.}
\bien{C}{defiende una propuesta («deberían recoger el agua») con razones.}{Tiene tesis y argumentos: gratuita, menos presión
sobre acueductos, menos inundaciones.}
\mal{D}{usa verbos en pasado.}{No los usa; eso sería propio del \S\,1.}
\end{comentada}

\taller
\piso{Calienta}{clasificar}
\begin{ejercicios}
\item Di la tipología de cada fragmento (\S\,1 a \S\,4). \sol{1 narrativo; 2 descriptivo; 3 expositivo; 4 argumentativo.}
\item Señala en el \S\,2 dos marcas de descripción. \sol{Por ejemplo: la comparación «fina como polvo» y «un brillo de vidrio molido»; los verbos de estado.}
\end{ejercicios}
\piso{Como en el examen}{opción múltiple}
\begin{ejercicios}
\item El \S\,3 tiene como propósito principal
\opcionesl{convencer al lector de cuidar el agua.}{explicar cómo se forma la lluvia.}{contar una experiencia personal.}{describir un paisaje.}
\sol{B.}
\item En el \S\,1, la expresión que marca el paso del tiempo es
\opcionesd{«empapado y feliz»}{«A las cinco»}{«en bicicleta»}{«tocó la puerta»}
\sol{B.}
\end{ejercicios}

\begin{retoresuelto}
El fragmento da órdenes en imperativo («disuelva», «revuelva», «consérvelo») para que el lector haga algo paso a paso.
\textbf{Opción C.}
\end{retoresuelto}
\bitacora

% ------------------------------------------------------------------
\leccion{pu-tt-2}{Géneros del examen}{Juego de emparejar cada texto con su género.}

\begin{reto}
Un texto presenta un libro recién publicado, resume su contenido sin contar el final y termina recomendándolo con algunas
reservas. Es
\opcionesd{una crónica}{una reseña}{un cuento}{una noticia}
\end{reto}

\tecnica{La técnica: el género dice qué esperar}
Dentro de cada tipología hay \clave{géneros} con reglas propias. Los que más aparecen en el examen:

\begin{center}\small
\begin{tabularx}{\linewidth}{@{}l>{\raggedright\arraybackslash}X@{}}
\toprule
\textbf{Género} & \textbf{Qué es y qué buscar}\\\midrule
Cuento & narración breve de ficción; buscar el conflicto y el cambio del personaje\\
Crónica & relato de hechos reales con mirada de autor; mezcla datos, escenas y voces\\
Noticia & hechos recientes, de lo más a lo menos importante; responde qué, quién, cuándo, dónde\\
Ensayo & reflexión personal y argumentada sobre un tema; buscar la tesis\\
Artículo de divulgación & explica un saber especializado a un público general\\
Reseña & presenta y valora una obra (libro, película, disco)\\
Editorial & opinión del medio sobre un hecho de actualidad\\
Carta & texto dirigido a un destinatario, con propósito explícito\\
Entrevista & preguntas y respuestas; distinguir la voz del entrevistador y la del entrevistado\\\bottomrule
\end{tabularx}
\end{center}

\begin{regla}
\begin{itemize}
  \item ¿Ficción o hechos reales? (cuento frente a crónica).
  \item ¿Informa o valora? (noticia o divulgación frente a reseña, editorial o ensayo).
  \item ¿A quién se dirige? (carta: destinatario explícito).
\end{itemize}
\end{regla}

\begin{texto}{«La casa de las once puertas», de Lucía Arango}{Reseña · texto original escrito para este libro (la novela es imaginaria)}
\pa Hay novelas que se leen de un tirón y otras que exigen detenerse. \emph{La casa de las once puertas}, primera novela de
Lucía Arango, pertenece a las segundas, y no por dificultad sino por densidad: cada capítulo es una puerta que se abre a
una época distinta de la misma casa, en un pueblo del suroeste antioqueño, a lo largo de un siglo.

\pa La estructura es su mayor acierto. En lugar de seguir a una familia, la novela sigue al edificio: lo vemos construirse
con los ahorros de un arriero, convertirse en escuela, en cuartel, en pensión, en ruina y, al final, en biblioteca. Los
personajes entran y salen, pero las paredes recuerdan.

\pa No todo funciona igual de bien. Los capítulos centrales, dedicados a los años de violencia, son tan contenidos que a
veces parecen apurados, como si la autora hubiera temido detenerse en el dolor. Y el lenguaje, por momentos, se enamora de
sí mismo: hay metáforas tan elaboradas que distraen de lo que cuentan.

\pa Aun así, es un debut notable. Arango consigue que el lector termine mirando de otro modo las casas viejas de su propio
pueblo, preguntándose qué habrán visto. Es una lectura recomendable para quien no tenga prisa.
\end{texto}

\subsection{Preguntas comentadas}
\begin{comentada}{¿Cuál es la valoración general que hace el autor de la reseña?}
\mal{A}{La novela es un fracaso por su lenguaje.}{\emph{Dice de más}: el lenguaje es una crítica parcial («por
momentos»).}
\bien{B}{Es un buen debut, con algunos defectos.}{\S\,4: «un debut notable»; \S\,3: lo que «no funciona igual de
bien».}
\mal{C}{Es una novela perfecta.}{\emph{Dice lo contrario}: el \S\,3 señala defectos.}
\mal{D}{No se puede saber qué opina.}{La opinión es explícita.}
\end{comentada}

\taller
\piso{Calienta}{reconocer}
\begin{ejercicios}
\item ¿Qué partes típicas de una reseña encuentras en el texto? \sol{Presentación de la obra y la autora (\S\,1), descripción de la estructura (\S\,2), crítica de debilidades (\S\,3) y valoración final con recomendación (\S\,4).}
\end{ejercicios}
\piso{Como en el examen}{opción múltiple}
\begin{ejercicios}
\item Según la reseña, el mayor acierto de la novela es
\opcionesd{su lenguaje}{los capítulos de la violencia}{su estructura}{su final}
\sol{C. \S\,2.}
\item La frase «los personajes entran y salen, pero las paredes recuerdan» \pf{2} significa que
\opcionesl{la casa es un fantasma.}{la casa es el verdadero hilo de la novela.}{los personajes no son importantes para la autora.}{la casa se destruye.}
\sol{B.}
\item Al decir que la novela es para «quien no tenga prisa», el autor
\opcionesl{advierte que exige una lectura pausada.}{critica que sea muy larga.}{recomienda leerla en vacaciones.}{dice que es aburrida.}
\sol{A. Coherente con el \S\,1 («exigen detenerse»).}
\end{ejercicios}

\begin{retoresuelto}
Presentar una obra, resumirla y valorarla es lo propio de la reseña. \textbf{Opción B.}
\end{retoresuelto}
\bitacora

% ------------------------------------------------------------------
\leccion{pu-tt-3}{Estructura y testimonio}{Juego de ordenar las partes de un texto.}

\begin{reto}
Un texto empieza describiendo cómo se perdían las cosechas de un pueblo por las heladas, sigue con la idea de un grupo de
campesinos de cubrir los cultivos con plástico de noche y termina con los resultados del primer año. Su estructura es
\opcionesd{cronológica pura}{problema y solución}{comparación}{enumeración}
\end{reto}

\tecnica{La técnica: reconocer el esqueleto}
Además de su tipo, todo texto tiene una \clave{estructura}: la forma en que se ordenan sus partes. Las más comunes:
\begin{itemize}
  \item \clave{Cronológica}: los hechos en el orden en que ocurrieron.
  \item \clave{Problema-solución}: se plantea una dificultad y luego una o varias respuestas.
  \item \clave{Causa-efecto}: se explica por qué ocurre algo o qué consecuencias tiene.
  \item \clave{Comparación-contraste}: se ponen dos cosas lado a lado.
  \item \clave{Tesis-argumentos}: una postura y sus razones (en los argumentativos).
\end{itemize}

\subsubsection{El testimonio}
Un \clave{testimonio} es un relato en primera persona de alguien que vivió los hechos. Su fuerza está en la experiencia
directa; su límite, en que es una sola mirada. En las preguntas, conviene distinguir lo que el narrador \emph{vivió}, lo
que \emph{siente} y lo que \emph{opina}.

\figura{estructuras}{Cuatro esqueletos frecuentes de un texto.}

\begin{regla}
\begin{itemize}
  \item Identifica el esqueleto con los conectores: «primero, luego» (cronológica); «el problema es…, una salida»
  (problema-solución); «por eso, como consecuencia» (causa-efecto); «en cambio, mientras que» (contraste).
  \item En un testimonio, separa hechos vividos, sentimientos y opiniones.
\end{itemize}
\end{regla}

\begin{texto}{«Aprendí a llegar a pie»}{Testimonio · texto original escrito para este libro (la narradora es ficticia)}
\pa Cuando me nombraron enfermera del puesto de salud de la vereda, pensé que el problema iba a ser la falta de medicinas.
Me equivoqué: el problema era el camino. El puesto quedaba a dos horas de la carretera, y la mitad de las familias vivía a
otras dos horas de allí, monte arriba.

\pa Los primeros meses esperaba a que la gente bajara. Casi nadie bajaba. Las madres no podían dejar a los otros hijos
solos; los viejos no aguantaban la subida de regreso; los hombres trabajaban de sol a sol. Las vacunas se vencían en la
nevera mientras yo leía revistas viejas.

\pa Entonces cambié de idea: si la gente no podía venir, yo iría. Organicé un recorrido por las fincas, dos días a la
semana, con una mula prestada que cargaba la nevera portátil. Al principio me recibían con desconfianza; al tercer mes ya
me tenían café listo y una lista de vecinos que también querían verme.

\pa No sé si hice lo correcto en todo. Hubo semanas en que el puesto quedó cerrado y alguien bajó y no me encontró. Pero
cuando me fui, cinco años después, casi todos los niños de la vereda tenían el esquema de vacunas completo. Creo que la
salud rural se mide menos en pastillas que en kilómetros.
\end{texto}

\subsection{Preguntas comentadas}
\begin{comentada}{La estructura del texto es principalmente}
\mal{A}{comparación entre la ciudad y el campo.}{\emph{No lo dice}: no hay comparación con la ciudad.}
\bien{B}{problema (la gente no llegaba) y solución (los recorridos), con sus resultados.}{\S\,1–2 el problema, \S\,3 la
solución, \S\,4 el balance.}
\mal{C}{causa y efecto de las enfermedades rurales.}{\emph{Dice otra cosa}: no explica enfermedades.}
\mal{D}{una lista de consejos para enfermeras.}{No es un texto instructivo.}
\end{comentada}

\begin{comentada}{¿Cuál de estas afirmaciones es una \emph{opinión} de la narradora y no un hecho que vivió?}
\mal{A}{Las vacunas se vencían en la nevera.}{Hecho \pf{2}.}
\mal{B}{Organizó recorridos dos días a la semana.}{Hecho \pf{3}.}
\bien{C}{La salud rural se mide menos en pastillas que en kilómetros.}{Es su conclusión personal \pf{4}, introducida con
«Creo que».}
\mal{D}{Casi todos los niños quedaron con el esquema completo.}{Hecho que ella informa \pf{4}.}
\end{comentada}

\taller
\piso{Como en el examen}{opción múltiple}
\begin{ejercicios}
\item Según el texto, la narradora se equivocó al pensar que el problema principal sería
\opcionesd{el camino}{la falta de medicinas}{la desconfianza}{la mula}
\sol{B. \S\,1.}
\item En el \S\,4, la narradora reconoce que
\opcionesl{su método tuvo costos, como dejar el puesto cerrado algunas semanas.}{su trabajo fue un fracaso.}{nunca debió aceptar el cargo.}{las familias no la querían.}
\sol{A.}
\item El cambio en la actitud de las familias se muestra en que
\opcionesl{bajaban más al puesto.}{al tercer mes le tenían café y una lista de vecinos.}{le prestaron una mula.}{pidieron otra enfermera.}
\sol{B. \S\,3.}
\end{ejercicios}

\begin{retoresuelto}
Se plantea una dificultad (las heladas), una respuesta (cubrir los cultivos) y sus resultados. \textbf{Opción B.}
\end{retoresuelto}
\bitacora

% ------------------------------------------------------------------
\evaluacion
\begin{texto}{Carta a quien hereda mi bicicleta}{Carta · texto original escrito para este libro}
\pa Querida Juana: te escribo porque mañana me voy a estudiar a Bogotá y la bicicleta se queda contigo. Antes de que te la
lleves, quiero contarte algunas cosas que no dice el manual.

\pa El freno de atrás chilla cuando llueve; no es que esté dañado, es que tiene carácter. La cadena se sale si cambias de
piñón en plena subida: cambia antes, en lo plano, como quien se prepara para una mala noticia. Y el asiento está bajo
porque así lo dejé el día que aprendí a montar sin rueditas; súbelo, que ya creciste.

\pa Esta bicicleta me llevó al colegio cuatro años. Con ella aprendí que las subidas se hacen despacio y que las bajadas
dan miedo solo al principio. También aprendí a arreglar un neumático pinchado en la mitad de la carretera, con las manos
negras y un vecino mirando.

\pa Cuídala, pero no tanto: una bicicleta guardada se oxida más rápido que una que se usa. Y cuando también te toque
irte, regálasela a alguien más pequeño. Así la bicicleta seguirá siendo de la familia aunque cambie de dueño.

\pa Te quiere, tu hermana, Laura.
\end{texto}
\begin{ejercicios}
\item El texto es principalmente
\opcionesd{un instructivo técnico}{una carta personal}{una crónica}{un ensayo}
\sol{B. Destinataria explícita, despedida y firma. (2.2)}
\item El \S\,2 combina
\opcionesl{instrucciones prácticas con un tono afectuoso y humorístico.}{datos científicos sobre bicicletas.}{una queja contra el fabricante.}{una narración de un accidente.}
\sol{A. (2.1)}
\item La expresión «tiene carácter» \pf{2}, referida al freno, sugiere que
\opcionesl{está dañado.}{tiene un comportamiento propio que no es una falla.}{es nuevo.}{hay que cambiarlo.}
\sol{B. «No es que esté dañado». (2.1)}
\item ¿Cuál de estas afirmaciones es una enseñanza que Laura dice haber aprendido con la bicicleta?
\opcionesl{Las subidas se hacen despacio.}{Las bajadas nunca dan miedo.}{Hay que guardar la bicicleta para que no se oxide.}{Los vecinos siempre ayudan.}
\sol{A. \S\,3. (2.3)}
\item El consejo de regalar la bicicleta a alguien más pequeño \pf{4} expresa el deseo de que
\opcionesl{Juana no la use mucho.}{la bicicleta conserve un vínculo familiar.}{se venda bien.}{nadie más la monte.}
\sol{B. (2.3)}
\end{ejercicios}
